\exercise{Barbell graph}{3}
Consider a barbell graph
\begin{center}
\includegraphics[width=0.3\textwidth]{barbel}
\end{center}
Construct the Laplacian $\bf L$ compute its eigenvalues and eigenvectors. Identify the algebraic connectivity $\lambda_2$ and compare it to the Cheeger's constant $g$ and identify the cut implied by the Fiedler vector $\vec{v_2}$. 

\solution 
Let's start by identifying Cheeger's constant for the graph. Intuitively the best way to separate the barbell is to cut it in the middle, which we can write as (1 2 3 | 4 5 6). This partition cuts one link and separates three nodes from the rest. So Cheeger's constant is 
\eq{
h= \frac13
}
Let's see if we can find this cut and constant also algebraically. The Adjacency matrix is 
\eq{
{\bf A} = \avecccccc{ 
  0 & 1 & 1 & 0 & 0 & 0 \\
  1 & 0 & 1 & 0 & 0 & 0 \\ 
  1 & 1 & 0 & 1 & 0 & 0 \\
  0 & 0 & 1 & 0 & 1 & 1 \\
  0 & 0 & 0 & 1 & 0 & 1 \\
  0 & 0 & 0 & 1 & 1 & 0 }
}
which gives us the degree matrix
\eq{
{\bf K} = \avecccccc{ 
  2 & 0 & 0 & 0 & 0 & 0 \\
  0 & 2 & 0 & 0 & 0 & 0 \\ 
  0 & 0 & 3 & 0 & 0 & 0 \\
  0 & 0 & 0 & 3 & 0 & 0 \\
  0 & 0 & 0 & 0 & 2 & 0 \\
  0 & 0 & 0 & 0 & 0 & 2 }
}
and hence the Laplacian 
\eq{
{\bf L} = \avecccccc{ 
   2 & -1 & -1 &  0 &  0 &  0 \\
  -1 &  2 & -1 &  0 &  0 &  0 \\ 
  -1 & -1 &  3 & -1 &  0 &  0 \\
   0 &  0 & -1 &  3 & -1 & -1 \\
   0 &  0 &  0 & -1 &  2 & -1 \\
   0 &  0 &  0 & -1 & -1 &  2 }
}
Next we have to find the eigenvalues and eigenvectors of the matrix. 
Straight away we know that there must be one eigenvalue $\lambda_1=0$ because the network has one component. The corresponding eigenvector is
\eq{
\vec{v_1}=(1,1,1,1,1,1)^{\rm T}
}
Also our symmetry tricks from the previous chapter should work for any matrix so based on the symmetry we node 1 and node 2 we can guess the eigenvector 
\eq{
\vec{v_{\rm x}}=(1,-1,0,0,0,0){\rm T} 
}
I have given it the number $x$ we don't quite know yet how large its eigenvalue is compared to the others. The first step to finding out is to compute $\lambda_x$. For this purpose I multiply the eigenvector with the Laplacian. As the eigenvector only lives on one side of the barbell, we can save some space by only multiplying the top $3\times 3$ block. So, we multiply 
\eq{
 \aveccc{ 
   2 & -1 & -1  \\
  -1 &  2 & -1  \\ 
  -1 & -1 &  3  }
\avec{1 \\ -1 \\ 0 } = \avec{3 \\ -3 \\ 0}
}
which shows that this is indeed an eigenvector and the corresponding eigenvalue is
\eq{\lambda_x=3}
We also have a corresponding eigenvalue on the other side. So, by analogy. 
\eq{
\vec{v_{\rm y}}=(0,0,0,0,1,-1){\rm T}, 
}
\eq{
\lambda_y = 3.
}
All other eigenvectors must have the same values in the pair symmetric nodes that are on either end of the barbell, but we are not done with symmetry yet. We can also exploit the horizontal symmetry of the barbell, so there should be eigenvectors of the form $(a,a,b,-b,-a,-a)^{\rm T}$. Let's try   
\eq{
{\bf L} = \avecccccc{ 
   2 & -1 & -1 &  0 &  0 &  0 \\
  -1 &  2 & -1 &  0 &  0 &  0 \\ 
  -1 & -1 &  3 & -1 &  0 &  0 \\
   0 &  0 & -1 &  3 & -1 & -1 \\
   0 &  0 &  0 & -1 &  2 & -1 \\
   0 &  0 &  0 & -1 & -1 &  2 }
\avec{a \\ a \\ b \\ -b \\ -a \\ -a}
= 
\avec{a-b \\ a-b \\ 4b-2a \\ 2a-4b \\ b-a \\ b-a}
= \lambda \avec{a \\ a \\ b \\ -b \\ -a \\ -a}
}
The matrix gives us two independent conditions 
\eqa{
    \lambda a &=& a-b \\
    \lambda b &=& 4b-2a 
}
Neither $a=0$ or $b=0$ can be a solution so we can set $a=1$ solving the quadratic polynomial 

Given the system of equations derived from the eigenvalue problem $Lv = \lambda v$ for an eigenvector of the form $v = (a, a, b, -b, -a, -a)^T$:
\begin{align*}
    \lambda a &= a - b \\
    \lambda b &= -2a + 4b
\end{align*}
We set $a=1$ and solve for $b$ and $\lambda$.

%\section*{Step 1: Substitute a = 1}

Substituting $a=1$ into the system gives two equations with two unknowns, $b$ and $\lambda$:
\begin{align*}
    \lambda &= 1 - b \\
    \lambda b &= -2 + 4b
\end{align*}
We substitute the expression for $\lambda$ into the second equation 
\eqa{
    (1 - b)b &=& -2 + 4b \\
    b - b^2 &=& -2 + 4b \\
    0 &=& b^2 + 3b - 2
}
We solve with the quadratic formula, which yields
\eq{
    b = \frac{-3 \pm \sqrt{3^2 - 4(1)(-2)}}{(2)(1)} = \frac{-3 \pm \sqrt{17}}{2} 
}
This gives us two solutions the first one is 
\eq{
  \lambda_z = 1-b = 1+\frac{3 +\sqrt{17}}{2} \approx 4.562
}
\eq{
  \vec{v_z} = (1,1,(-3-\sqrt{17})/2,(3+\sqrt{17})/2,-1,-1)^{\rm T}
}
The other is 
\eq{
\lambda_2 = 1-b = 1+\frac{3 -\sqrt{17}}{2} \approx 0.438
}
\eq{
\vec{v_2} = (1,1,(3+\sqrt{17})/2,(-3-\sqrt{17})/2,-1,-1)^{\rm T}
}
where I have guessed that this will probably the Fiedler vector because its shape is consistent with the partition we expect and the eigenvalue is so low. 
To be sure we need to compute the final eigenvalue that we are still missinng, and show that it is greater than $\lambda_2$. 

A clever way to find the last eigenvalue is to note that the trace of the Laplacian is twice the sum over all node degrees, 
\eq{
{\rm Tr}({\bf L}) = 2+2+3+3+2+2 = 14 
}
but also it must be the sum over eigenvalues
\eqa{
{\rm Tr}({\bf L}) &=& 0 + \frac{5 -\sqrt{17}}{2} + \frac{5 +\sqrt{17}}{2} +3 + 3 + \lambda_m \\
   &=& 11 + \lambda_m
}
where $\lambda_m=3$ is the missing eigenvalue. Of course if we don't trust this reasoning we can always check it. Because we have found all the odd-symmetry eigenvector it must have even symmetry, so we try 
\eq{
{\bf L} = \avecccccc{ 
   2 & -1 & -1 &  0 &  0 &  0 \\
  -1 &  2 & -1 &  0 &  0 &  0 \\ 
  -1 & -1 &  3 & -1 &  0 &  0 \\
   0 &  0 & -1 &  3 & -1 & -1 \\
   0 &  0 &  0 & -1 &  2 & -1 \\
   0 &  0 &  0 & -1 & -1 &  2 }
\avec{a \\ a \\ b \\ b \\ a \\ a}
= 
\avec{a-b \\ a-b \\ 2b-2a \\ 2b-2a \\ a-b \\ a-b}
}
which yields the two conditions
\eqa{
\lambda a &=& a-b \\
\lambda b &=& 2b-2a 
}
One solution is $a=b$ and $\lambda=0$, but we already knew this. But mutiplying the the first row with 2 and adding the second gives us 
\eq{
  \lambda (2a+b) = 0 
}
Since we are not interested in $\lambda=0$ we can divide by $\lambda$ which yields 
\eq{
2a+b=0
}
There fore $a=1$, $b=-2$ and $\lambda=3$ is a solution, so we can confirm 
\eq{
\lambda_m=3 
}
\eq{
\vec{v_m} = (1,1,-2,-2,1,1)^{\rm T}
}
In summary we have found the eigenvalues 
\eq{
  0, 0.438 , 3, 3, 3, 4.562
}
Moreover we have shown that the Fiedler vector is 
\eq{
\vec{v_2} = (1,1,(3+\sqrt{17})/2,(-3-\sqrt{17})/2,-1,-1)^{\rm T}
}
which confirms that the optimal solution to the ratio-cut problem is given by the partition (1 2 3 | 4 5 6).  

As expected the algebraic connectivity $\lambda_2 \approx  0.438$ is similar to the Cheeger's constant $h\approx 0.333$. Moreover, we can check that the results obey Cheeger's inequality
\eq{
\frac{\lambda_2}{2k_{\rm max}} \leq h \leq \sqrt{2k_{\rm max}\lambda_2}
}
We compute the bounds
\eq{
\frac{\lambda_2}{2k_{\rm max}} = \frac{\lambda_2}{6} \approx 0.073 
}
\eq{
\sqrt{2k_{\rm max}\lambda_2} = \sqrt{6\lambda_2} \approx 1.621
}
So the inequality holds, 
\eq{
0.073 \leq 0.333 \leq 1.621, 
}
but the bounds are not very tight. 
